\pagestyle{empty}
\tight
\begin{tblr}{width=\textwidth,colspec={|Q[c,m,1.9cm]|X[l,m]|},hlines,vlines,hspan=minimal,row{1}={bg=headblue},rowsep=1pt,colsep=3pt}
\SetCell{c,m}\hfil\logo\hfil & \textbf{POLITEKNIK NEGERI MALANG}\par\textbf{JURUSAN TEKNOLOGI INFORMASI}\par\textbf{PROGRAM STUDI: D4 TEKNIK INFORMATIKA} \\
\SetCell[c=2]{c} \textbf{RENCANA PEMBELAJARAN SEMESTER (RPS)} & \\
\end{tblr}
\begin{longtblr}{width=\textwidth,colspec={|Q[l,m,1.9cm]|X[0.85,l,m]|X[1.45,l,m]|X[0.75,l,m]|X[1.15,l,m]|},hlines,vlines,hspan=minimal,rowsep=1pt,colsep=2pt,row{1,3}={bg=headgray,font=\bfseries}}
MATA KULIAH & KODE & BOBOT (SKS)/jam & SEMESTER & TGL. PENYUSUNAN \\
Metode Numerik & RTI253006 & 2 SKS/4 jam & 3 & 6 Juli 2026 \\
OTORISASI & Dosen Pengembang RPS & Koordinator RMK & \SetCell[c=2]{l} Ka PRODI &  \\
 & Adevian Fairuz Pratama, S.ST., M.Eng & Vivi Nur Wijayaningrum, S.Kom., M.Kom. & \SetCell[c=2]{l}{\vspace{8mm}\par Dr. Ely Setyo Astuti, S.T., M.T.} &  \\
\SetCell[r=11]{l} Capaian Pembelajaran (CP) & \SetCell[c=4]{l,bg=headgray,font=\bfseries} Capaian Pembelajaran Lulusan Program Studi (CPL-Prodi) &  &  &  \\
 & CPL09 & \SetCell[c=3]{l} Mampu memahami cara kerja sistem komputer dan menggunakan pendekatan komputasi untuk penyelesaian masalah secara efektif. &  &  \\
 & CPL10 & \SetCell[c=3]{l} Memiliki kompetensi untuk menganalisis persoalan kompleks untuk mengidentifikasi solusi bidang informatika/ilmu komputer dengan mempertimbangkan wawasan ilmu multidisiplin. &  &  \\
 & \SetCell[c=4]{l,bg=headgray,font=\bfseries} Capaian Pembelajaran mata kuliah (CPL-MK) &  &  &  \\
 & CPMK0902 & \SetCell[c=3]{l} Mampu menerapkan prinsip matematis dan pendekatan komputasi untuk menganalisis dan menyelesaikan persoalan komputasional secara logis dan sistematis. &  &  \\
 & CPMK1007 & \SetCell[c=3]{l} Mampu mengembangkan solusi numerik berbasis pendekatan komputasi diskrit untuk menyelesaikan persoalan kompleks bidang informatika dan teknik. &  &  \\
 & \SetCell[c=4]{l,bg=headgray,font=\bfseries} Kemampuan akhir tiap tahapan belajar (Sub-CPMK) &  &  &  \\
 & SCPMK0902-02401 & \SetCell[c=3]{l} Mahasiswa mampu menjelaskan konsep metode numerik, galat, dan menyelesaikan sistem persamaan linear menggunakan metode Gauss, Gauss-Jordan, dan Gauss-Seidel. &  &  \\
 & SCPMK0902-02402 & \SetCell[c=3]{l} Mahasiswa mampu menyelesaikan persamaan nonlinear menggunakan metode tertutup (biseksi, regula falsi) dan metode terbuka (Newton-Raphson, sekans). &  &  \\
 & SCPMK0902-02403 & \SetCell[c=3]{l} Mahasiswa mampu menyelesaikan masalah diferensiasi numerik dan integrasi numerik menggunakan metode Riemann, trapezoidal, dan Simpson. &  &  \\
 & SCPMK1007-02404 & \SetCell[c=3]{l} Mahasiswa mampu menerapkan interpolasi dan regresi numerik untuk aproksimasi fungsi dan analisis data. &  &  \\
\SetCell[r=2]{l} Penilaian & \SetCell[c=4]{l} *Nilai CPMK didapat dari agregasi nilai SUB-CPMK yang dibebankan &  &  &  \\
 & \SetCell[c=4]{l}{\tiny\begin{tblr}{width=\linewidth,colspec={|Q[c,m,0.1\linewidth]|Q[c,m,0.16\linewidth]|X[l,m]|X[l,m]|X[l,m]|X[l,m]|X[l,m]|Q[c,m,0.08\linewidth]|},hlines,vlines,hspan=minimal,rowsep=1pt,colsep=0.7pt,row{1,2}={bg=headgray,font=\bfseries}}
Kode CPMK & Kode SCPMK & \SetCell[c=5]{c} Bobot per Bentuk Penilaian & & & & & Total Bobot \\
 & & Tugas & Kuis 1 & UTS & Kuis 2 & UAS & \\
\SetCell[r=3]{c} CPMK0902 & SCPMK0902-02401 & 8\%: tugas galat dan persamaan linear & 10\%: konsep numerik dan galat & 10\%: persamaan linear & & & 28\% \\
 & SCPMK0902-02402 & 4\%: tugas persamaan nonlinear & & 15\%: persamaan nonlinear & & & 19\% \\
 & SCPMK0902-02403 & 5\%: tugas diferensiasi dan integrasi & & & 10\%: diferensiasi dan integrasi & 15\%: integrasi numerik & 30\% \\
\SetCell[r=1]{c} CPMK1007 & SCPMK1007-02404 & 3\%: tugas interpolasi dan regresi & & & & 20\%: interpolasi dan regresi & 23\% \\
\SetCell[c=7]{r}\textbf{Total Per Penilaian} & & & & & & & \textbf{100\%} \\
\end{tblr}} &  &  &  \\
Diskripsi Singkat Mata Kuliah & \SetCell[c=4]{l} Mata kuliah Metode Numerik membekali mahasiswa dengan kemampuan menyelesaikan masalah matematika teknik secara numerik menggunakan komputer. Topik meliputi konsep galat, penyelesaian sistem persamaan linear (Gauss, Gauss-Jordan, Gauss-Seidel), persamaan nonlinear (metode tertutup dan terbuka), diferensiasi numerik, integrasi numerik (Riemann, trapezoidal, Simpson), interpolasi, dan regresi. Implementasi dilakukan menggunakan perangkat lunak komputasi. &  &  &  \\
Bahan Kajian / Materi Pembelajaran & \SetCell[c=4]{l}\begin{enumerate}\item Pengantar metode numerik dan galat\item Sistem persamaan linear: Gauss, Gauss-Jordan, Gauss-Seidel\item Persamaan nonlinear: metode tertutup dan terbuka\item Diferensiasi numerik\item Integrasi numerik: Riemann, trapezoidal, Simpson\item Interpolasi dan regresi\end{enumerate} &  &  &  \\
\SetCell[r=2]{l} Pustaka & \SetCell[c=4]{l} \textbf{Utama:}\begin{enumerate}\item Rahmad, E.C., Ikawati, D.S.E., \& Syaifudin, Y.W. 2018. \textit{Metode Numerik}. Polinema Press. ISBN 9786026695888.\item Burden, R.L. \& Faires, J.D. 2011. \textit{Numerical Analysis}, 9th ed. Brooks/Cole Cengage Learning.\item Chapra, S.C. \& Canale, R.P. 2015. \textit{Numerical Methods for Engineers}, 7th ed. McGraw-Hill.\end{enumerate} &  &  &  \\
 & \SetCell[c=4]{l} \textbf{Pendukung:}\begin{enumerate}\item Heath, M.T. 2002. \textit{Scientific Computing: An Introductory Survey}, 2nd ed. McGraw-Hill.\item Kreyszig, E. 2011. \textit{Advanced Engineering Mathematics}, 10th ed. Wiley.\item Materi LMS, spreadsheet MS Excel, dan lembar kerja praktikum.\end{enumerate} &  &  &  \\
Media Pembelajaran & \SetCell[c=2]{l} \textbf{Software:} MS Excel, Python (NumPy/SciPy), MATLAB/Octave, LMS. &  & \SetCell[c=2]{l} \textbf{Hardware:} Komputer/laptop, proyektor, kalkulator ilmiah. &  \\
Nama Dosen Pengampu & \SetCell[c=4]{l}\begin{enumerate}\item Adevian Fairuz Pratama, S.ST., M.Eng\item Triana Fatmawati, S.T., M.T.\item Dian Hanifudin Subhi, S.Kom., M.Kom.\item M. Hasyim Ratsanjani, S.Kom., M.Kom.\end{enumerate} &  &  &  \\
Matakuliah Syarat & \SetCell[c=4]{l} RTI251004 Matematika Dasar &  &  &  \\
\end{longtblr}
\begin{longtblr}{width=\textwidth,colspec={|Q[c,m,0.06\textwidth]|X[1.35,l,m]|X[1.08,l,m]|X[1.16,l,m]|X[0.7,l,m]|X[1.24,l,m]|X[1.06,l,m]|X[0.98,l,m]|Q[c,m,0.05\textwidth]|},hlines,vlines,hspan=minimal,rowhead=2,row{1,2}={bg=headgreen,font=\bfseries},rowsep=1pt,colsep=1.5pt}
Mgg.\par Ke & Kemampuan Akhir Yang Direncanakan (Sub-CP-MK) & Bahan kajian (Materi Pembelajaran) & Bentuk dan Metode Pembelajaran & Estimasi Waktu & Pengalaman Belajar Mahasiswa & Kriteria \& Bentuk Penilaian & Indikator Penilaian & Bobot Penilaian (\%) \\
(1) & (2) & (3) & (4) & (5) & (6) & (7) & (8) & (9) \\
1 & SCPMK0902-02401 & Pengantar Metode Numerik: pengertian metode numerik vs analitik, motivasi dan aplikasi, sumber galat, representasi bilangan floating point. & Modalitas: Blended Learning. Bentuk: Luring/Daring. Metode: ceramah, tanya jawab, latihan soal, dan penyelesaian studi kasus. & 1 x 2 x 50' tatap muka; 1 x 2 x 50' tugas/praktik mandiri. & Mahasiswa membedakan metode analitik dan numerik, mengidentifikasi sumber galat, dan menghitung representasi bilangan floating point pada contoh kasus. & Bentuk penilaian: Tugas pengantar metode numerik. Mengacu rubrik RTM. & Ketepatan membedakan metode analitik dan numerik; identifikasi sumber galat. & 2 \\
2 & SCPMK0902-02401 & Galat (Error Analysis): galat pemotongan (truncation), galat pembulatan (rounding), galat relatif dan absolut, propagasi galat, dan signifikansi angka besar. & Modalitas: Blended Learning. Bentuk: Luring/Daring. Metode: ceramah, tanya jawab, latihan soal, dan penyelesaian studi kasus. & 1 x 2 x 50' tatap muka; 1 x 2 x 50' tugas/praktik mandiri. & Mahasiswa menghitung galat pemotongan, galat pembulatan, dan galat relatif pada contoh perhitungan numerik, serta menganalisis propagasi galat pada operasi aritmetika. & Bentuk penilaian: Tugas analisis galat. Mengacu rubrik RTM. & Ketepatan perhitungan galat; kemampuan menganalisis propagasi galat. & 2 \\
3 & SCPMK0902-02401 & Sistem Persamaan Linear -- Eliminasi Gauss: bentuk baku SPL, eliminasi maju (forward elimination), substitusi mundur (back substitution), pivoting parsial, dan deteksi singularitas. & Modalitas: Blended Learning. Bentuk: Luring/Daring. Metode: ceramah, tanya jawab, latihan soal, dan penyelesaian studi kasus. & 1 x 2 x 50' tatap muka; 1 x 2 x 50' tugas/praktik mandiri. & Mahasiswa menyelesaikan sistem persamaan linear menggunakan eliminasi Gauss dengan pivoting parsial dan memverifikasi solusi secara manual maupun menggunakan perangkat lunak. & Bentuk penilaian: Tugas eliminasi Gauss. Mengacu rubrik RTM. & Ketepatan proses eliminasi Gauss; verifikasi solusi dan penanganan pivoting. & 2 \\
4 & SCPMK0902-02401 & Sistem Persamaan Linear -- Gauss-Jordan dan Gauss-Seidel: eliminasi Gauss-Jordan (reduksi baris penuh), metode iteratif Gauss-Seidel, kriteria konvergensi, dan perbandingan efisiensi. & Modalitas: Blended Learning. Bentuk: Luring/Daring. Metode: ceramah, tanya jawab, latihan soal, dan penyelesaian studi kasus. & 1 x 2 x 50' tatap muka; 1 x 2 x 50' tugas/praktik mandiri. & Mahasiswa menyelesaikan SPL menggunakan Gauss-Jordan dan iterasi Gauss-Seidel, memeriksa konvergensi, serta membandingkan efisiensi kedua metode. & Bentuk penilaian: Tugas Gauss-Jordan dan Gauss-Seidel. Mengacu rubrik RTM. & Ketepatan iterasi Gauss-Seidel; pemeriksaan konvergensi dan perbandingan metode. & 2 \\
5 & SCPMK0902-02401 & Kuis 1: materi minggu 1--4 (pengantar metode numerik, analisis galat, eliminasi Gauss, Gauss-Jordan, Gauss-Seidel). & Modalitas: Luring. Bentuk: Ujian tertulis, sifat closed book. & Sesuai jadwal asesmen. & Mahasiswa mengerjakan soal kuis teori dan perhitungan manual metode numerik secara mandiri. & Bentuk penilaian: Kuis 1 konsep numerik dan galat. Mengacu rubrik RTM. & Ketepatan jawaban konsep dan perhitungan; kemampuan menerapkan metode pada kasus. & 10 \\
6 & SCPMK0902-02402 & Persamaan Nonlinear -- Metode Tertutup: metode biseksi (bolzano), metode regula falsi (false position), teorema nilai antara, kriteria berhenti, dan analisis galat iterasi. & Modalitas: Blended Learning. Bentuk: Luring/Daring. Metode: ceramah, tanya jawab, latihan soal, dan penyelesaian studi kasus. & 1 x 2 x 50' tatap muka; 1 x 2 x 50' tugas/praktik mandiri. & Mahasiswa menyelesaikan persamaan nonlinear menggunakan biseksi dan regula falsi, melacak tabel iterasi, serta menganalisis konvergensi dan galat setiap iterasi. & Bentuk penilaian: Tugas metode tertutup. Mengacu rubrik RTM. & Ketepatan tabel iterasi; analisis galat dan konvergensi metode biseksi dan regula falsi. & 2 \\
7 & SCPMK0902-02402 & Persamaan Nonlinear -- Metode Terbuka: metode titik tetap (fixed point), Newton-Raphson, metode sekans, perbandingan laju konvergensi, dan analisis kondisi divergensi. & Modalitas: Blended Learning. Bentuk: Luring/Daring. Metode: ceramah, tanya jawab, latihan soal, dan penyelesaian studi kasus. & 1 x 2 x 50' tatap muka; 1 x 2 x 50' tugas/praktik mandiri. & Mahasiswa menyelesaikan persamaan nonlinear menggunakan Newton-Raphson dan sekans, membandingkan laju konvergensinya dengan metode tertutup, serta mengidentifikasi kondisi divergensi. & Bentuk penilaian: Tugas metode terbuka. Mengacu rubrik RTM. & Ketepatan iterasi Newton-Raphson dan sekans; analisis laju konvergensi dan kondisi divergensi. & 2 \\
8 & SCPMK0902-02402 & Persamaan Nonlinear -- Aplikasi dan Studi Kasus: penerapan metode numerik untuk mencari akar persamaan pada kasus teknik dan informatika, implementasi dengan perangkat lunak. & Modalitas: Blended Learning. Bentuk: Luring/Daring. Metode: ceramah, tanya jawab, latihan soal, dan penyelesaian studi kasus. & 1 x 2 x 50' tatap muka; 1 x 2 x 50' tugas/praktik mandiri. & Mahasiswa menerapkan metode pencarian akar persamaan nonlinear pada kasus nyata teknik/informatika dan mengimplementasikannya menggunakan perangkat lunak komputasi. & Bentuk penilaian: Tugas studi kasus persamaan nonlinear. Mengacu rubrik RTM. & Ketepatan penerapan metode pada kasus nyata; kualitas implementasi perangkat lunak. & 0 \\
9 & SCPMK0902-02401, SCPMK0902-02402 & UTS: materi minggu 1--8 (galat, persamaan linear, persamaan nonlinear). & Modalitas: Luring. Bentuk: Ujian tertulis, sifat closed book. & Sesuai jadwal asesmen. & Mahasiswa mengerjakan soal UTS teori dan perhitungan metode numerik secara mandiri. & Bentuk penilaian: UTS metode penyelesaian persamaan. Mengacu rubrik RTM. & Ketepatan perhitungan; kemampuan analisis dan pemilihan metode yang tepat. & 25 \\
10 & SCPMK0902-02403 & Diferensiasi Numerik: deret Taylor, diferensiasi maju (forward difference), mundur (backward difference), tengah (central difference), tingkat akurasi, dan galat pemotongan. & Modalitas: Blended Learning. Bentuk: Luring/Daring. Metode: ceramah, tanya jawab, latihan soal, dan penyelesaian studi kasus. & 1 x 2 x 50' tatap muka; 1 x 2 x 50' tugas/praktik mandiri. & Mahasiswa menghitung turunan fungsi secara numerik menggunakan diferensiasi maju, mundur, dan tengah, kemudian menganalisis akurasi dan galat pemotongan masing-masing metode. & Bentuk penilaian: Tugas diferensiasi numerik. Mengacu rubrik RTM. & Ketepatan perhitungan turunan numerik; analisis akurasi dan galat pemotongan. & 2 \\
11 & SCPMK0902-02403 & Integrasi Numerik -- Metode Riemann: integral Riemann (kiri, kanan, tengah), hubungan dengan luas area, estimasi galat integrasi, dan penerapan pada fungsi sederhana. & Modalitas: Blended Learning. Bentuk: Luring/Daring. Metode: ceramah, tanya jawab, latihan soal, dan penyelesaian studi kasus. & 1 x 2 x 50' tatap muka; 1 x 2 x 50' tugas/praktik mandiri. & Mahasiswa menghitung integral fungsi menggunakan metode Riemann kiri, kanan, dan tengah, membandingkan hasilnya dengan nilai eksak, dan menganalisis galat integrasi. & Bentuk penilaian: Tugas integrasi Riemann. Mengacu rubrik RTM. & Ketepatan perhitungan integral Riemann; perbandingan dengan nilai eksak dan analisis galat. & 1 \\
12 & SCPMK0902-02403 & Kuis 2: materi minggu 10--11 (diferensiasi numerik dan integrasi Riemann). & Modalitas: Luring. Bentuk: Ujian tertulis, sifat closed book. & Sesuai jadwal asesmen. & Mahasiswa mengerjakan soal kuis teori dan perhitungan diferensiasi dan integrasi numerik secara mandiri. & Bentuk penilaian: Kuis 2 diferensiasi dan integrasi numerik. Mengacu rubrik RTM. & Ketepatan perhitungan turunan dan integral numerik; analisis galat. & 10 \\
13 & SCPMK0902-02403 & Integrasi Numerik -- Metode Trapezoidal: aturan trapezoid tunggal, trapezoid komposit, estimasi galat, dan perbandingan akurasi dengan Riemann. & Modalitas: Blended Learning. Bentuk: Luring/Daring. Metode: ceramah, tanya jawab, latihan soal, dan penyelesaian studi kasus. & 1 x 2 x 50' tatap muka; 1 x 2 x 50' tugas/praktik mandiri. & Mahasiswa menghitung integral menggunakan aturan trapezoidal tunggal dan komposit, memverifikasi hasil dengan nilai eksak, dan membandingkan akurasi dengan metode Riemann. & Bentuk penilaian: Tugas integrasi trapezoidal. Mengacu rubrik RTM. & Ketepatan perhitungan trapezoidal komposit; perbandingan akurasi antar metode. & 1 \\
14 & SCPMK0902-02403 & Integrasi Numerik -- Metode Simpson: aturan Simpson 1/3 dan 3/8, Simpson komposit, estimasi galat, penerapan pada kasus teknik, dan perbandingan metode integrasi numerik. & Modalitas: Blended Learning. Bentuk: Luring/Daring. Metode: ceramah, tanya jawab, latihan soal, dan penyelesaian studi kasus. & 1 x 2 x 50' tatap muka; 1 x 2 x 50' tugas/praktik mandiri. & Mahasiswa menghitung integral menggunakan aturan Simpson 1/3 dan 3/8 komposit, membandingkan akurasi dengan trapezoidal dan Riemann, serta menerapkannya pada kasus teknik. & Bentuk penilaian: Tugas integrasi Simpson. Mengacu rubrik RTM. & Ketepatan perhitungan Simpson komposit; perbandingan akurasi dan penerapan pada kasus. & 1 \\
15 & SCPMK1007-02404 & Interpolasi: interpolasi linear, polinomial Lagrange, polinomial Newton (beda terbagi), pemilihan titik interpolasi, dan galat interpolasi. & Modalitas: Blended Learning. Bentuk: Luring/Daring. Metode: ceramah, tanya jawab, latihan soal, dan penyelesaian studi kasus. & 1 x 2 x 50' tatap muka; 1 x 2 x 50' tugas/praktik mandiri. & Mahasiswa menghitung nilai aproksimasi fungsi menggunakan interpolasi Lagrange dan Newton, membandingkan akurasi berbagai derajat polinomial, serta menganalisis galat interpolasi. & Bentuk penilaian: Tugas interpolasi. Mengacu rubrik RTM. & Ketepatan perhitungan interpolasi Lagrange dan Newton; analisis galat dan pemilihan derajat. & 1.5 \\
16 & SCPMK1007-02404 & Regresi: regresi linear sederhana (least squares), regresi linear berganda, regresi nonlinear (linearisasi), koefisien determinasi $R^2$, dan penerapan pada analisis data. & Modalitas: Blended Learning. Bentuk: Luring/Daring. Metode: ceramah, tanya jawab, latihan soal, dan penyelesaian studi kasus. & 1 x 2 x 50' tatap muka; 1 x 2 x 50' tugas/praktik mandiri. & Mahasiswa menentukan persamaan regresi linear dan nonlinear dari data eksperimen menggunakan metode kuadrat terkecil, menghitung $R^2$, dan menginterpretasikan hasilnya. & Bentuk penilaian: Tugas regresi. Mengacu rubrik RTM. & Ketepatan perhitungan regresi linear; interpretasi koefisien determinasi dan kualitas model. & 1.5 \\
17 & SCPMK0902-02403, SCPMK1007-02404 & UAS: materi minggu 1--16 dengan penekanan pada integrasi numerik, interpolasi, dan regresi. & Modalitas: Luring. Bentuk: Ujian tertulis, sifat closed book. & Sesuai jadwal asesmen. & Mahasiswa mengerjakan soal UAS teori dan perhitungan metode numerik komprehensif secara mandiri. & Bentuk penilaian: UAS implementasi metode numerik lengkap. Mengacu rubrik RTM. & Ketepatan perhitungan; kemampuan analisis dan pemilihan metode yang tepat. & 35 \\
\end{longtblr}
